% 光的干涉（高中）
% keys 波的叠加
% license Usr
% type Tutor

\subsection{双缝干涉}
\subsubsection{理论推导}
\subsubsection{实验装置}
一个经典的双缝干涉实验装置如下图所示。在单缝前面有一滤光片，用以得到单色光（角频率相同）。在把光波视作球面波后，我们可以画出一系列同心圆，每一个圆代表一个相位，因此，到达双缝的总是同一相位，也就是光源的初始相位。由此，我们在双缝处得到两个相位差为0的光源。
\begin{figure}[ht]
\centering
\includegraphics[width=12cm]{./figures/17f92705aa639da5.png}
\caption{} \label{fig_interf_1}
\end{figure}

\subsection{薄膜干涉}